\section{Common mass, smoothing, and score moments}\label{sec:tools}
The induction needs three operations: converting overlap into hidden offsets, controlling multilinear score expansions, and conditioning on an event where the reference density is bounded below.  We prove them with the uniformity needed later.

\begin{lemma}[Multilinear moment bound]\label{lem:moment}
Let $t\ge2$ be an even integer, let $h\ge0$, and let $(Z_j)_{j\in\mathcal J}$ be a finite family of independent centered real variables with $\|Z_j\|_t\le h$.  For an integer $r\ge1$, the degree-$r$ multilinear
polynomial with real coefficients $A(S)$ on $r$-subsets satisfies
\begin{equation}\label{eq:hc}
\left\|\sum_{|S|=r} A(S)\prod_{j\in S}Z_j\right\|_t
\le (C_1 t h)^r \|A\|_{\ell^2},
\end{equation}
where $\ell^2$ uses counting measure and $C_1$ is absolute.
\end{lemma}

\begin{proof}
The tensorization follows the classical principle underlying hypercontractive estimates \citep{Bonami1970}.  We give an elementary argument that applies to arbitrary independent centered real coordinates satisfying the stated moment bound; it requires neither symmetry nor a common distribution.  Normalize $h=1$ (the zero case is clear). A scalar satisfies
$\|b+vZ\|_t^2\le |b|^2+C_2 t |v|^2$ for real $b,v$:
if $|v|\le |b|/t$ and $b\ne0$, expand the moment, with the linear
term zero, to bound it by $|b|^t(1+O(t^2(v/b)^2))$
using $\sum_{k\ge2}(t|v/b|)^k/k!$. Here concavity gives $(1+O(t^2(v/b)^2))^{2/t}\le
1+O(t(v/b)^2)$. Otherwise use the triangle inequality and $2|bv|\le 2t |v|^2$. Iteration by conditioning
on one variable's complement and applying the triangle inequality
in $L^{t/2}$ proves \eqref{eq:hc}. We will also apply this to finite
families indexed by positions/ranks, identifying ordered increasing
tuples with subsets.

\end{proof}

\begin{lemma}[Lifting adjacent-power overlap]\label{lem:lift}
Let $w$ be a nonempty word, $k$ a positive integer, and $0<\omega<1/8$ such that $\ov(\Tr(w^k),\Tr(w^{k+1}))\ge\omega$.  There is an absolute constant $C_{\rm abs}$ such that, whenever $C\ge C_{\rm abs}/\omega$, $P$ is a sufficiently large positive integer, and the integer $F$ satisfies $F\ge CkP$, we have
\begin{equation}\label{eq:lift}
\ov(\Tr(w^{FP}),\Tr(w^{FP+P}))\ge 2/3.
\end{equation}
It is enough that $P$ exceed an absolute constant, uniformly in $w,k,\omega,C,F$.
\end{lemma}

\begin{proof}
This is a quantitative common-part smoothing statement, related in mechanism to consecutive-power smoothing for positive contractions \citep[Propositions 2.2--2.3]{Zaharopol1989}.  Write the words as a common tail of $FP-g(k+1/2)$ copies
and $g=2\lfloor FP/(2k+1)\rfloor$ chunks of $k$ or $k+1$ copies.
Use a uniform random schedule conditional on the total number
$m$ of long chunks being $g/2$ or $g/2+P$ respectively.
Let $\omega U$ be a common sublaw of chunk trace distributions. Use
independent flags of probability $\omega$ to emit the trace under $U$
(hiding the label $k$ or $k+1$), otherwise reveal the label and
generate from the normalized residual for that label. The augmented
laws postprocess to the claimed ordinary ones by joining the chunks
and simulating the known tail. Given the flags and the sequence of
visible labels, the other data have the same conditional law on both
hypotheses. The number $M$ of hidden flags belongs to
$[\omega g/2,2\omega g]$ except with probability at most
$2\exp(-c'\omega g)$ for an absolute $c'>0$ by the multiplicative
binomial Chernoff bound, where
$\omega g\ge c'' \omega C P^2$ for an absolute $c''>0$ at all large $P$.
Conditional on such flags, the total variation between visible-label
sequences equals that of visible long counts $v$ since the sequences
are exchangeable on the visible indices. For total $m$ this count
has hypergeometric law $p_m$, with
\[
 p_m(v)={\binom{g-M}{v}\binom{M}{m-v}}/{\binom{g}{m}}.
\]
Uniformly for $g/2\le m\le g/2+P$,
$\sup_v p_m(v)=O((\omega g)^{-1/2})$. To see this, consecutive
probability ratios locate a mode within $O(1)$ of $(g-M)m/g$,
e.g.\ a mode is $\lfloor(g-M+1)(m+1)/(g+2)\rfloor$. For large $P$
the counts and complements in the numerator at the mode are all
fixed positive fractions (say between $1/4$ and $3/4$)
of the respective trial totals because $g\gtrsim C P^2$,
$M\ge\omega g/2\gtrsim P^2$ and $g-M\ge (1-2\omega)g$.
Stirling then gives the bound (the exponential
entropy factor is at most one by concavity). Couple adjacent $m,m+1$
by uniformly flipping one short label. This increases $v$ by one
with probability $t(v)=(g-M-v)/(g-m)$ conditional on $v$.
Hence $p_{m+1}(v)-p_m(v)=t(v-1)p_m(v-1)-t(v)p_m(v)$. The sequence
$t(v)p_m(v)$ on the support, extended by zero, is unimodal by
log-concavity. Thus the adjacent total variation is $\max_v t(v)p_m(v)$,
bounded by $\max_v p_m(v)$. The full variation cost of all $P$ steps
is $O(P/\sqrt{\omega g})$ plus the rare flag probability, at most
$1/3$ at large $P$ by choosing $C\omega$ greater than a sufficiently
large absolute constant. Taking smaller common sublaws when necessary lets
us use exactly mass $\omega$ in the flag construction.
\end{proof}

\subsection{Paired hiding and the reported information}
The paired construction below uses a joint common component of mass $\omega^2$; on this component, the orientation remains fair conditional on the flag and both reported side outputs.  Common-part constructions appear in the average-case block argument of \citet[Section 3.3]{MPV2014} and the partial-trace construction in the proof of \citet[Lemma 2.1]{HoldenLyons2020}.

Figure~\ref{fig:paired-hiding} depicts one pair and the resulting uncertainty in the number of copies to the left of the central segment.
\begin{figure}[htbp]
\centering
\begin{tikzpicture}[x=1cm,y=1cm,font=\small,>=Latex,
  side/.style={draw=linkblue,fill=linkblue!6,minimum width=2.8cm,
               minimum height=0.48cm,inner sep=2pt},
  center/.style={draw=linkblue,fill=linkblue!18,minimum width=3.5cm,
                 minimum height=0.48cm,inner sep=2pt},
  card/.style={draw=linkblue,rounded corners=2pt,align=center,
               text width=5.8cm,minimum height=0.94cm,inner sep=5pt}]
  \node[anchor=west,font=\small\bfseries] at (0,0.58)
    {A fair bit swaps the side lengths};
  \node[anchor=east] at (13.7,0.58)
    {$E\sim\mathrm{Bernoulli}(1/2)$};

  \node[anchor=east] at (0.75,0) {$E=0$};
  \node[side] at (2.4,0) {$w^k$};
  \node[center] at (7.0,0) {central segment};
  \node[side] at (11.6,0) {$w^{k+1}$};
  \draw[linkblue] (3.8,0)--(5.25,0);
  \draw[linkblue] (8.75,0)--(10.2,0);

  \node[anchor=east] at (0.75,-0.65) {$E=1$};
  \node[side] at (2.4,-0.65) {$w^{k+1}$};
  \node[center] at (7.0,-0.65) {central segment};
  \node[side] at (11.6,-0.65) {$w^k$};
  \draw[linkblue] (3.8,-0.65)--(5.25,-0.65);
  \draw[linkblue] (8.75,-0.65)--(10.2,-0.65);

  \node[draw=linkblue,fill=white,rounded corners=2pt,
        minimum width=12.6cm,minimum height=0.54cm] (law) at (7.0,-1.48)
    {$P_E=\omega^2(U\otimes U)+(1-\omega^2)R_E$};
  \node[card,fill=linkblue!9] (hidden) at (3.45,-2.52)
    {\textbf{Hidden:} \(E\) is not reported.\\[2pt]
     Emit \(U\otimes U\); \(E\) remains fair.};
  \node[card,draw=black!50,fill=black!3] (visible) at (10.55,-2.52)
    {\textbf{Visible:} report \(E\).\\[2pt]
     Emit the joint residual \(R_E\).};
  \draw[->,linkblue] (4.8,-1.75)--(3.45,-1.99);
  \draw[->,black!65] (9.2,-1.75)--(10.55,-1.99);

  \node[draw=linkblue,fill=linkblue!5,rounded corners=2pt,
        text width=13.0cm,align=center,inner sep=6pt,anchor=north] at (7.0,-3.50)
    {\(J\) pairs contribute
     \(L_{\rm left}:=Jk+\sum_{i=1}^{J}E_i\) copies of \(w\) on the left.\\[3pt]
     After the side report, \(c_{\rm vis}:=Jk+\sum_{\rm visible}E_i\)
     and the hidden count \(M\) are fixed, and\\[3pt]
     \(\displaystyle
       L_{\rm left}=c_{\rm vis}+\sum_{\rm hidden}E_i
          \ \overset{\mathrm{law}}{=}\ 
       c_{\rm vis}+\mathrm{Bin}(M,1/2).\)};
\end{tikzpicture}
\caption{Let \(\omega U\) be a common sublaw of the trace laws of
\(w^k\) and \(w^{k+1}\). For orientation \(E\), \(P_E\) is the product
of its two side laws and \(R_E\) is the normalized joint residual.
Across \(J\) independent pairs straddling the segment,
\(M\sim\mathrm{Bin}(J,\omega^2)\) before reporting. The lower panel
conditions on flags, side outputs, and visible bits; \(L_{\rm left}\)
counts only copies contributed by the left pair members. Other fixed
material adds a deterministic location offset. Block widths are schematic.}
\label{fig:paired-hiding}
\end{figure}


Throughout, unless window boundaries or additional payloads are
specified explicitly, report only one unseparated output
per segment/chunk/residual, plus flags and revealed orientation bits
as indicated (not original word boundaries internal to that output).
Given a word $w$ and a count $k$ with $\ov(\Tr(w^k),\Tr(w^{k+1}))\ge\omega$, choose a common sublaw $\omega U$.  Two chunks use a paired random bit orientation giving lengths
$(k,k+1)$ or $(k+1,k)$ even when far separated in the concatenation.
Use one joint flag of probability $\omega^2$ and output from
$U\otimes U$ on it, otherwise reveal the bit and generate the
two traces jointly from the normalized residual after subtracting
$\omega^2 U\otimes U$ from the product law for that orientation.
Marginalizing flags yields the requisite conditionally independent
pair traces given orientations. Conditional on the augmented
pair data, the bit remains fair on the hidden event. Independent
pairs thus carry independent hidden fair bits conditional on this
part of the report. All auxiliary schedules/flags are redrawn
independently for each trace.

\newpage
\subsection{Conditioning to obtain a small likelihood score}
Separating a typical region from a small exceptional mass is also used in Hellinger estimates for trace reconstruction \citep[arXiv v2, Section 2]{Chase2021Lower}.  The following version keeps the full moment and simulation-error budgets.

\begin{lemma}[Clipping with a common event]\label{lem:clipping}
Fix a level $0\le i\le D$ and its parameters from Section~\ref{sec:scales}.  Suppose preliminary oracle laws have common
``side'' metadata with law $Q_{\rm side}$, and conditional densities
$G_x,G_y$ over a conditional reference $Q$, possibly depending on
metadata.  Require $\int G_x\,dQ=\int G_y\,dQ=1$ for $Q_{\rm side}$-almost every metadata value.  Use $Q_{\rm side}Q$ as the joint reference. Suppose
on a typical metadata event of probability $\ge1-o(\Delta)$ we have
uniform estimates
\begin{equation}\label{eq:clipcond}
\|G_x-1\|_{r_i},\|G_y-1\|_{r_i}\le N_i^{-\rho},
\qquad \|G_y-G_x\|_{r_i}^2\le N_i^{-d_i-.01}
\end{equation}
in $L^{r_i}(Q)$, where $\rho\ge 2^{-4B-10}$.
Then conditioning each law on the event consisting of typical metadata and $G_x\ge1/2$ changes it in total variation by $o(\Delta)$.  The conditioned laws $\mu,\nu$ satisfy $\nu\ll\mu$ and
\[
 \|d\nu/d\mu-1\|_{L^{r_i}(\mu)}^2\le N_i^{-d_i}.
\]
All these conclusions hold uniformly over $0\le i\le D$.
\end{lemma}

\begin{proof}
Condition each preliminary law separately on the same joint event
$A_{\mathrm{keep}}$ consisting of typical metadata and $G_x\ge1/2$. Its complement has joint reference
measure on the typical part $\le(2N_i^{-\rho})^{r_i}$, hence bad
mass $o(\Delta)$ under each preliminary law by Cauchy--Schwarz,
the moment bound, and the common atypical mass. The conditioned
laws therefore change postprocessing by $o(\Delta)$ in total variation.
Because $G_x\ge1/2$ on the conditioning event $A_{\mathrm{keep}}$,
the new $y$-law is dominated by the new $x$-law and their likelihood
ratio on the latter support is $c' G_y/G_x$ with a scalar
$c'=1+o(\Delta)$. Also
\[
 \int_{\{G_x\ge1/2,\ \mathrm{typical}\}}
  \left|G_y/G_x-1\right|^{r_i} G_x\,d(Q_{\rm side}Q)
 \le 2^{r_i-1} N_i^{-(d_i+.01)r_i/2}.
\]
Under the conditioned $x$-law this yields
$\|G_y/G_x-1\|_{r_i}\le 3 N_i^{-(d_i+.01)/2}$ for all large $n$
by taking the $r_i$-th root (the conditioning probability is $1-o(\Delta)$).
Thus
$\|c'G_y/G_x-1\|_{r_i}\le
c' \|G_y/G_x-1\|_{r_i}+|c'-1|\le N_i^{-d_i/2}$
as desired for large $n$, since $\Delta\le N_i^{-100 d_D}$.
In this calculation and throughout, any $o(\Delta)$ masses claimed
from \eqref{eq:clipcond} follow uniformly using the definition of
$r_i$ and $\log N_i\to\infty$; no uniform lower bound on the
unconditioned ratio or density is needed. For Hellinger comparisons
of preliminary laws one can similarly use the squared density
difference times $2$ on $G_x\ge1/2$ and $G_x+G_y$ on its complement.

\end{proof}
